problem:  
Let \(M^m \subset \mathbb{S}^{m+k}(1)\) be a compact, properly biharmonic submanifold with constant mean curvature vector field \(H\).  
The biharmonic equation can be stated as  
\[
\Delta^\perp H = mH - \operatorname{trace}B(A_H(\cdot),\cdot),
\]
and for such \(M\), the induced Laplacian identity  
\[
\frac{1}{2}\Delta |H|^2 = \|\nabla^\perp H\|^2 + \langle \operatorname{trace}B(A_H(\cdot),\cdot),H\rangle - m|H|^2
\]
shows how the normal derivative of the mean curvature interacts with the curvature tensor.  

In the minimal case, the analogous formula for the second fundamental form \(A\) led Simons to the classical identity  
\[
\frac{1}{2}\Delta |A|^2 = \|\nabla A\|^2 + m|A|^2 - \sum_{i,j}\langle A_i,A_j\rangle^2,
\]
which underpins the *gap theorems* for minimal submanifolds in spheres.  
A general formula of this type in the **biharmonic** setting is unknown.

---

**Research Problem (Biharmonic Simons Identity and Gap Phenomenon):**  
Derive a Simons-type identity for the squared norm of \(A_H\) or \(\nabla^\perp H\) in the setting of constant-mean-curvature properly biharmonic submanifolds of spheres; that is, find a formula of the form  
\[
\frac{1}{2}\Delta |A_H|^2 = \|\nabla A_H\|^2 + \Phi(|A_H|^2,|H|^2,\text{curvature terms}),
\]
where \(\Phi\) is an explicit algebraic polynomial in curvature quantities depending at most on \(m\) and \(k\).  

Once such a formula is found, investigate whether one can extract from it an *integral gap inequality* of the type  
\[
\int_M (\|A_H\|^2 - f(|H|^2))^2 dV \ge C_m \int_M \|\nabla A_H\|^2 dV,
\]
where \(f\) represents the model relation for the known proper biharmonic examples and \(C_m>0\) is dimensionally determined.  
In particular, determine whether a Simons-type structure implies that any deviation of \(|H|^2\) from the canonical value \(1/2\) forces a curvature‑energy gap in \(\|A_H\|^2\).

---

Why is it a "good" problem:  
- **Conceptual clarity and structural base:** Unlike arbitrary inequalities, this problem begins with the quest for a genuine *biharmonic Simons identity*, directly modeled on the most fundamental analytic tool in submanifold geometry. It is not speculative; it seeks to uncover a missing structural equation.  
- **Bridging classical and higher‑order geometries:** It extends the central ideas of minimal submanifold theory (Simons identity, gap theorems) to the fourth‑order regime of biharmonic immersions, blending intrinsic analysis and extrinsic curvature algebra.  
- **Analytic and geometric depth:** Producing and analyzing such an identity would require high‑order Weitzenböck formulas, refined curvature pinching, and intricate Bochner‑type computations. These analytic challenges go beyond what is known for harmonic maps.  
- **Simplicity of statement, profound implications:** It asks, essentially, “Does a Simons‑type identity and the ensuing curvature‑energy gap exist for biharmonic submanifolds?”—a question short enough to be stated in one sentence but deep enough to define an entire research program.  
- **Predictive and unifying power:** A positive resolution would reveal a new rigidity mechanism parallel to the classical case, potentially leading to dimension‑and‑codimension‑independent classification results and thereby advancing the general theory of biharmonic maps.  

This formulation is mathematically grounded, natural within the established geometric framework, and poised to stimulate significant theoretical progress.